% ============================================================
% Section 1: Introduction
% ============================================================
% Target length: ~1 column (500–700 words)
% Write early (Month 4) — refine as project matures
% Structure: problem → why existing solutions fail →
%            our approach → contributions → paper outline
% ============================================================

\section{Introduction}
\label{sec:introduction}

% ── Paragraph 1: The problem ────────────────────────────────
% Hook: browsers were designed for humans, not agents
% Establish the governance vacuum

The web browser was designed thirty years ago for human users
clicking links and reading pages.
In 2026, AI agents autonomously browse, fill forms, execute purchases,
and orchestrate complex multi-step workflows inside that same
architecture---yet browsers expose no formal mechanism for governing,
auditing, or limiting what an agent can do.
The result is a governance vacuum with real consequences.

% ── Paragraph 2: Real attacks (cite them all) ───────────────
% Establish urgency — this is not theoretical

In 2025, this vacuum was repeatedly exploited.
CometJacking~\cite{layerx_cometjacking_2025} exfiltrated emails and
calendar data from Perplexity Comet via hidden instructions.
Screenshot Injection~\cite{brave_screenshot_injection_2025} hijacked
Comet and Fellou agents using text invisible to the human user.
Tainted Memories~\cite{layerx_tainted_memories_2025} poisoned
ChatGPT Atlas's long-term memory via CSRF.
HashJack~\cite{cato_hashjack_2025} embedded malicious instructions
in URL fragments.
In each case, the browser granted the agent the same privileges as
the human user, with no capability boundary and no audit trail.

% ── Paragraph 3: Why existing defenses fail ─────────────────
% Model-level defenses are insufficient — cite OpenAI admission

Existing defenses operate at the model level: adversarial training,
classifiers, and red-teaming.
OpenAI acknowledged in December 2025 that prompt injection
``may never be fully solved'' at the model
level~\cite{openai_prompt_injection_2025}.
A concurrent study by researchers at OpenAI, Anthropic, and
DeepMind tested twelve prompt injection defenses and found all
fail against adaptive attacks~\cite{nasr2025}.
The conclusion is clear: model-level defenses are insufficient,
and architectural enforcement is necessary.

% ── Paragraph 4: Our approach ───────────────────────────────
% Introduce Ferrite, the key insight, and the architecture

We present \ferrite{}, a browser built from the ground up to
enforce capability-based governance over all agent and extension
actions.
\ferrite{}'s central insight is that the capability broker---not
the rendering engine, not the agent, not the extension---is the
sole privileged boundary.
Every action touching sensitive resources requires an unforgeable
capability token, evaluated against composable Rego policies,
and recorded in a tamper-evident cryptographic audit log.
A successfully prompt-injected agent cannot exfiltrate data if
the capability to read mail was never granted.
A malicious extension cannot perform unauthorized network requests
if its token has already expired.
The architecture enforces security independently of model behavior.

% ── Paragraph 5: Contributions ──────────────────────────────

\noindent This paper makes the following contributions:

\begin{itemize}
  \item \textbf{Architecture.} The first formal capability-based
    governance model for LLM browser agents, combining the
    Capsicum/COWL lineage~\cite{watson2010capsicum,stefan2014cowl}
    with Certificate Transparency-style audit
    logs~\cite{laurie2013rfc6962}.

  \item \textbf{Implementation.} \ferrite{}, a working browser
    prototype in Rust on the Servo engine, with a policy-evaluated
    capability broker, hash-chained audit log upgrading to a Merkle
    tree, and a Wasm extension sandbox with capability-gated host
    functions.

  \item \textbf{Developer tools.} An Extension Security Analyzer
    (CLI) and an Agent Action Inspector (UI), usable independently
    of the browser.

  \item \textbf{Evaluation.} Demonstration that \ferrite{} blocks
    \todo{N} adversarial scenarios including reproductions of the
    four 2025 real-world attacks, with median broker latency of
    \todo{X}$\mu$s---showing architectural enforcement is both
    effective and efficient.
\end{itemize}

% ── Paragraph 6: Paper outline ──────────────────────────────

The remainder of this paper is structured as follows.
\S\ref{sec:background} provides background on capability systems
and agentic browsers.
\S\ref{sec:threat_model} defines our threat model and adversary.
\S\ref{sec:design} presents \ferrite{}'s architecture and design.
\S\ref{sec:implementation} describes the implementation.
\S\ref{sec:evaluation} presents our evaluation.
\S\ref{sec:related_work} discusses related work.
\S\ref{sec:conclusion} concludes.
